\section{Method and System Architecture}\label{sec:method}

\subsection{Research method}
This study applies the Design Science Research (DSR) method \parencite{hevner2004design}, a methodological framework for information-systems research in which knowledge is produced through the deliberate construction and evaluation of technical artifacts. In DSR, an \emph{artifact} is a designed technical object that addresses a specific problem; here it is the entire Continux architecture comprising the K3s infrastructure, the streaming pipeline from Vector to Iceberg, the GitOps layer using Argo CD, and the observability layer based on VictoriaMetrics and Grafana. \emph{Evaluation} is the phase that examines the artifact through controlled experimental runs rather than merely describing intended configuration on paper.

Evaluation evidence consists of command output as text files, metric responses returned by VictoriaMetrics in JSON form, and Grafana dashboard captures. Output produced on the deployment host (node \texttt{imac}) was copied to a Windows analysis workspace at \path{evidence/finalize/20260522-151720/} so that figures and numbers could be cross-checked while writing the report. Only state, counts, and timings present in evidence are reported as outcomes; properties for which no corresponding measurement exists -- including end-to-end exactly-once reconciliation and multi-load benchmarks defined in the repository -- are listed explicitly as limitations rather than inferred.

\subsection{Dataset and replay preparation}
The New York City Taxi and Limousine Commission (NYC TLC) publishes public taxi trip records, one of the most commonly used datasets in urban-mobility analytics research \parencite{nyctlc}. Continux uses Yellow Taxi data for \texttt{2026-03}. A converter script transforms the source Parquet file -- a compressed columnar format optimized for batch analytics -- into JSONL to produce an ordered event stream for replay. The TLC Taxi Zone lookup table contains \texttt{265} rows and is kept as a CSV on MinIO; each row maps a \texttt{location\_id} to a \texttt{borough} and a \texttt{zone}.

\begin{table*}[t!]
\centering
\caption{Experimental data source}
\begin{tabularx}{0.86\textwidth}{@{}p{0.38\textwidth}Y@{}}
\toprule
\textbf{Property} & \textbf{Value}\\
\midrule
Trip type and month & Yellow Taxi, \texttt{2026-03}\\
Source/output formats & Parquet / JSONL\\
Confirmed JSONL rows & \texttt{3,952,451}\\
Lookup data & TLC Taxi Zone, \texttt{265} rows\\
Reported run & Controlled replay, final public MV at \texttt{986} trips\\
\bottomrule
\end{tabularx}
\end{table*}

Each JSONL row is a serialized taxi trip; when Vector emits the row into the broker, it becomes an event in the replay. Vector adds two meta fields on emission: \texttt{event\_id} (a unique identifier for each emitted event) and \texttt{event\_time} (the wall-clock moment of emission). The deployed manifest rate-limits replay to \texttt{2} events per second so that the resource-constrained lab cluster can process the stream stably; Vector is deliberately stopped after the observation target is reached. Therefore the reported MV is not a full-file processing result; a deeper analysis of raw-data quality is provided in Section~\ref{sec:dataquality}.

\subsection{Overall architecture}
\begin{figure*}[t!]
\centering
\includegraphics[width=0.65\textwidth]{architecture-process-flow.png}
\caption{Continux architecture and process flow.}
\label{fig:architecture-en}
\end{figure*}

Figure~\ref{fig:architecture-en} presents three concurrent paths. In the data path, Parquet is converted to JSONL for event replay, Redpanda buffers the stream, RisingWave maintains SQL results, and the sink writes Iceberg objects on MinIO. In the control path, GitHub holds manifests while Argo CD deploys and reconciles Kubernetes resources. In the observation path, workloads and the exporter expose metrics for VictoriaMetrics storage and Grafana visualization.

Keeping these paths distinct clarifies the measurement boundary: Iceberg objects establish that the data output path ran; Argo CD status establishes the state of core deployed configuration; query-loop and cutover metrics establish whether the public query name remained usable during the logic switch.

\subsection{Deployment topology}
\begin{table*}[t!]
\centering
\caption{Experimental K3s cluster topology}
\label{tab:topology-en}
\begin{tabularx}{\textwidth}{@{}p{0.18\textwidth}p{0.40\textwidth}Y@{}}
\toprule
\textbf{Node} & \textbf{Recorded resources} & \textbf{Role}\\
\midrule
\texttt{imac} & Intel i5-8500, 6 cores, 8 GB RAM, 200 GB SSD, Ubuntu 26.04 & Server \#1; Vector, Redpanda, RisingWave, MinIO, exporter\\
\texttt{continux-vps} & 2 vCPU, 4 GB RAM, 80 GB SSD, Ubuntu 24.04 & Server \#2; Argo CD, VictoriaMetrics, Grafana\\
\texttt{helios-pc} & WSL Ubuntu 26.04, Intel i5-12500H host, 16 GB RAM & Server \#3; quorum-only, node-exporter\\
\bottomrule
\end{tabularx}
\end{table*}

Nodes communicate through \emph{Tailscale} -- a mesh VPN built on the WireGuard protocol that creates an overlay network among hosts that may be in different physical networks, as if they shared a LAN. K3s in HA mode uses \emph{embedded etcd}: \emph{etcd} is a distributed key-value store that holds the Kubernetes cluster state, and ``embedded'' means it runs inside the K3s process rather than separately. The cluster requires a \texttt{2/3} \emph{quorum}: at least two of three etcd members must remain online for write operations to reach consensus, which is why \texttt{helios-pc} is reserved as a quorum-only voter and does not run data workloads. Compute and storage workloads are placed on \texttt{imac} to exploit its local SSD and RAM; \texttt{continux-vps} hosts the control plane and observability stack.

\clearpage
\onecolumn
\subsection{SQL objects and Blue/Green design}
\begin{table}[H]
\centering
\caption{Core data objects}
\begin{tabularx}{\textwidth}{@{}p{0.27\textwidth}p{0.27\textwidth}Y@{}}
\toprule
\textbf{Object} & \textbf{Input} & \textbf{Function}\\
\midrule
\texttt{tlc\_zone} & CSV on MinIO & Location-to-zone lookup\\
\texttt{nyc\_taxi\_src} & Redpanda topic & JSON event source\\
\texttt{mv\_zone\_stats\_blue} & Source + lookup & Baseline aggregate\\
\texttt{mv\_zone\_stats} & Initial blue MV & Public queried name\\
\texttt{mv\_zone\_stats\_green} & Source + lookup + filter & Candidate logic\\
\texttt{sink\_zone\_stats} & Public MV at sink creation & Iceberg upsert output\\
\bottomrule
\end{tabularx}
\end{table}

The public name \texttt{mv\_zone\_stats} is the client query contract: clients need not modify SQL when serving logic changes from blue to green. Conversely, \texttt{sink\_zone\_stats} was created from a dependency that existed before the swap, so green results at the public MV do not by themselves establish that the Iceberg sink automatically changed logic.

\begin{lstlisting}[language=SQL,caption={Additional filtering in the green logic}]
WHERE t.fare_amount >= 0
  AND t.trip_distance >= 0
\end{lstlisting}

Once the green MV was populated, a loop queried the public MV at a short interval while the system performed:

\begin{lstlisting}[language=SQL,caption={Blue/Green MV cutover command}]
ALTER MATERIALIZED VIEW mv_zone_stats
SWAP WITH mv_zone_stats_green;
\end{lstlisting}

\subsection{Evaluation measures}
\begin{table}[H]
\centering
\caption{Evaluation measure groups and interpretation}
\begin{tabularx}{\textwidth}{@{}p{0.23\textwidth}p{0.35\textwidth}Y@{}}
\toprule
\textbf{Group} & \textbf{Evidence} & \textbf{Interpretation}\\
\midrule
Cutover & Duration, query errors, restarts, counts & Public MV switched without observed query error\\
Observed integrity & MV counts, filter difference, objects & Output exists; difference is explained by logic\\
Streaming activity & Replay progress, dashboard, displayed lag & Demonstrates activity, not general benchmark\\
Resource stability & Nodes, workloads, PVC, CPU/RAM & Cluster sustains measured scenario\\
\bottomrule
\end{tabularx}
\end{table}

\subsection{Evidence provenance and traceability}
The local evidence bundle for the measured run is located at \path{evidence/finalize/20260522-151720/} and contains \texttt{48} textual or JSON output files. Baseline files record node, PVC, Argo CD, and RisingWave readiness; the \texttt{05-*} group records replay and Iceberg listing results; and the \texttt{06-*} group records green catch-up, the query loop, swap, duration, and metrics stored by VictoriaMetrics.

Evidence interpretation separates event time from later metric retrieval time. The swap timestamp recorded in \path{06-cutover-duration.txt} is epoch \texttt{1779466691}, corresponding to \texttt{2026-05-22 16:18:11 UTC}, within the query-loop observation window. VictoriaMetrics files confirm duration and query-error metrics at scrape timestamp \texttt{1779467656}, later than the swap; that timestamp establishes persisted observation rather than replacing the cutover time. Metadata in the source bundle shows that the swap/duration files were updated after the query-loop log and post-swap result; the duration is therefore reported as a runbook-recorded measure, not as latency independently reconstructed from an immutable trace.
